
\section{\trGenus{} $3$ 的\trCurve}

\textbf{19.7.1. \trCanonical{}嵌入。}

我们先就\trGenus{} $g>1$ 作一些一般性的讨论，然后再专门讨论\trGenus{} $3$。

\textbf{19.7.A.（重要练习）}\quad 设 $k=\bar{k}$（纯粹是为了避免分心——你尽可以去掉这条假设）。设 $C$ 是一条\trGenus{}为 $g$ 的\trCurve，而 $g>1$。证明：若 $C$ 不是超椭圆的，则\trCanonical{}\trInvertibleSheaf{}给出一个\trClosedEmbedding{} $C\hookrightarrow\mathbb{P}^{g-1}$。（在超椭圆的情形，我们已经见过：\trCanonical{}层给出的是到一个\trRationalNormalCurve{}的二重覆盖。）提示：回想 $h^0(C,\omega_C)=g$（\mbox{练习~19.2.A}）。使用我们在 \S19.1 中的那个「闭嵌入判别法」，即以其 \S19.2.10 的形式。回想：若 $h^0(C,\OO(p+q))=2$，则 $C$ 是超椭圆的（\mbox{练习~19.5.B}）。

这道练习给出了「一切\trGenus{} $2$ \trCurve{}都是超椭圆的」的第二个证明。

\textbf{19.7.2.}\ \emph{\trhDefinition{}。}\quad 这样一条经\trCanonical{}层（又称\trCanonical{}\trLineBundle）嵌入 $\mathbb{P}^{g-1}$ 的\trCurve{}称为\textbf{\trCanonical{}曲线}（canonical curve），而这个\trClosedEmbedding{}则称为 $C$ 的\textbf{\trCanonical{}嵌入}（canonical embedding）。

\textbf{19.7.3. 回到\trGenus{} $3$。}

现在设 $C$ 是一条\trGenus{} $3$ 的\trCurve。那么 $\omega_C$ 的\trDegree{}为 $2g-2=4$，并且带有 $g=3$ 个\trSection。由\mbox{练习~19.7.A}，当（且仅当）$C$ 不是超椭圆的时，\trCanonical{}映射便把 $C$ 描述为 $\mathbb{P}^2$ 中一条\trDegree{}为 $4$ 的\trCurve。

反过来，任何\trPlaneQuartic{}都是被典范地嵌入的。理由：这条\trCurve{}具有\trGenus{} $3$（见 \eqref{eq:genus-plane-curve}），并且由一个\trDegree{}为 $4$、带有三个\trSection{}的\trInvertibleSheaf{}所映射。但由\mbox{练习~19.2.A}，\trDegree{}为 $2g-2$ 且带有 $g$ 个\trSection{}的\trInvertibleSheaf{}只有 $\omega_C$。

特别地，每条非超椭圆的\trGenus{} $3$ \trCurve{}都可以恰好以一种方式描述为\trPlaneQuartic{}（在 $\mathbb{P}^2$ 的\trAutomorphism{}之下；这些\trAutomorphism{}已在\mbox{练习~15.5.A} 中得到描述）。我们已然证明：在模掉\trProjectiveTransformation{}（\S15.5.2）的意义下，非超椭圆的\trGenus{} $3$ \trCurve{}与\trPlaneQuartic{}之间存在一个\trBijection。

\textbf{19.7.4.}\ \emph{\trhRemark{}。}\quad 特别地，由于存在\trRegular{}的\trPlaneQuartic（\mbox{练习~13.3.E}），故存在非超椭圆的\trGenus{} $3$ \trCurve。

\textbf{19.7.B.（练习，参看 \S19.5.8）}\quad 给出一个启发式（非严格）的论证：非超椭圆的\trGenus{} $3$ \trCurve{}构成一个维数为 $6$ 的族。（提示：数一数\trRegular{}\trPlaneQuartic{}之族的维数，再模掉 $\Aut\mathbb{P}^2=\mathrm{PGL}(3)$。）

于是，\trGenus{} $3$ 的\trCurve{}看似分属两个族：超椭圆\trCurve{}（一个维数为 $5$ 的族）与非超椭圆\trCurve{}（一个维数为 $6$ 的族）。这是一种假象——它们实际上同属一个维数为 $6$ 的族。事实上，超椭圆\trCurve{}自然地就是非超椭圆\trCurve{}的极限。我们可以写下一个显式的族。（这番解释不可避免地需要一些「挥手而过」的成分\footnote{译注：原文 \emph{hand-waving}，指不作严格证明、以手势示意带过的论证方式。}，因为它牵涉到我们尚未见过的论题。）设我们有一条在 $\mathbb{P}^1$ 的 $2g+2=8$ 个点处分支的超椭圆\trCurve。选取一个把 $\mathbb{P}^1$ 与 $\mathbb{P}^2$ 中一条\trConic{}等同起来的\trIsomorphism。于是存在一条\trRegular{}的\trPlaneQuartic，恰好在那八个点处与这条\trConic{}相交。（这需要一个用到\trBertini{}定理~\mbox{13.4.2} 的简短论证，我们略去。）于是，若 $f$ 是这条\trConic{}的方程，而 $g$ 是这条四次曲线的方程，则 $f^2+t^2g$ 便是一族对大多数 $t$ 都\trSmooth{}的四次曲线（光滑性是 $t$ 的一个开条件，见关于一般光滑性的\mbox{定理~21.6.6}）。$t=0$ 的情形是一条二重\trConic。而事实是：若你把这个族的\trTotalSpace{}\trNormalization，则中心纤维（在 $t=0$ 之上者）就变成了我们的超椭圆\trCurve。这样，我们便把这条超椭圆\trCurve{}表示成了非超椭圆\trCurve{}的极限。

\textbf{19.7.C.（不重要练习）}\quad 把容许一个到 $\mathbb{P}^1$ 的\trDegree{} $3$ 覆盖的（射影）\trCurve{}（在域 $k$ 上）称为\textbf{\trTrigonal{}的}（trigonal）。证明：每条非超椭圆的\trGenus{} $3$ 复\trCurve{}都是\trTrigonal{}的：取其在 $\mathbb{P}^2$ 中的四次模型，并从曲线上任意一点向 $\mathbb{P}^1$ 作\trProjection。具体做法：在 $\mathbb{P}^2$ 上选取\trCoordinate{}，使得 $p$ 落在 $[0,0,1]$ 处。（同样的想法——施于三次曲线而非四次曲线——将在 \S19.9.9 中得到使用。）

\textbf{19.7.5.}\ \hard\ \textbf{一条没有非平凡\trAutomorphism{}的\trGenus{} $3$ \trCurve。}

我们已经看到：\trGenus{}至多为 $2$ 的（\trSmooth、\trProjective、\trIntegral{}的）\trCurve{}总有非平凡\trAutomorphism。而事实表明，存在没有非平凡\trAutomorphism{}的\trGenus{} $3$ \trCurve。

\textbf{19.7.D.（练习）}\quad 设 $C'\subset\mathbb{P}^2$ 是一条\trSmooth{}的\trPlaneQuartic{}（在任意域 $k$ 上）。证明：$C'$ 的\trAutomorphism{}与保持 $C'$ 的 $\mathbb{P}^2$ 的\trAutomorphism{}之间存在\trBijection\footnote{原书此处作「there is bijection between」，当脱一「a」，径补。——译者注}。

于是，要找一条没有非平凡\trAutomorphism{}的\trGenus{} $3$ \trCurve，我们只须找到一条\trSmooth{}的\trPlaneQuartic{} $C'$，使得 $\mathbb{P}^2$ 中把 $C'$ 作为集合予以保持的\trAutomorphism{}必为恒等映射。你的直觉或许会（正确地）告诉你：大多数四次曲线都具有这条性质。但要具体展示一个 $C'$（连同证明），就得挽起袖子真刀真枪地干了。\trPoonen{}在 [P1] 中给出了在任意域上都没有非平凡\trAutomorphism{}的\trCurve；一个特征 $0$ 的例子是
\[
y^3z-3yz^3=3x^4-4x^3z+z^4.
\]

\textbf{19.7.E.（练习）}\quad 设 $C$ 是一条没有非平凡\trAutomorphism{}的\trSmooth、\trProjective{}的\trCurve。证明：$C$ 的任何两个互不相同的\trOpenSubset{}都互不\trIsomorphic。

\textbf{19.7.6.}\ \emph{具有非平凡\trAutomorphism{}的\trGenus{} $3$ \trCurve。}

当然，\trGenus{} $3$ 的\trCurve{}可以有非平凡\trAutomorphism：超椭圆\trCurve{}便是明证。更令人赞叹的是\trKlein{}四次曲线
\[
x^3y+y^3z+z^3x=0
\]
（在 $\mathbb{P}^2_{\mathbb{C}}$ 中），它具有 $168$ 个\trAutomorphism。（你能把它们全都找出来吗？）事实上，\trKlein{}四次曲线的\trAutomorphismGroup{}是阶为 $168$ 的那个唯一的有限单群（即第二小的非阿贝尔有限单群）。在 \S21.4.9 中，我们将看到任何\trGenus{} $3$ \trCurve{}都不可能具有更多的\trAutomorphism（此处假定「\trGenus{}大于 $1$ 的\trCurve{}具有有限\trAutomorphismGroup」这一艰深的事实）。
